\scanpage{20}
\mspage{19}

\[
 \mathrm{Omb}(F)=\{X\in\unsure{M|\mathcal{M}}\mid X\ll F\}\,\deleted{\unsure{\subset|=}}
 \tag*{$(*)$}
\]

et la multiombre
\[
 \mathrm{Multom}(F)=\{\mathrm{Omb}(X)\mid X\in\widetilde F\}.
 \tag*{$(**)$}
\]
\margin{$\widetilde F=F$, si on identifi $F$ à une partie de $M$ via $\unsure{\mathcal{F}|\mathfrak{F}|\mathscr{F}}\unsure{\mapsto|\to}\unsure{\mathcal{P}|\mathrm{P}}(M)$\\ $F\unsure{\mapsto|\vdash}\widetilde F$}

En apparence $(*)$ utilise la connaissance de $\ll$ dans $\mathcal{F}$ tout entier, mais \deleted{\illegible{1}}\inserted{il} \deleted{\unsure{on dira qu'elle|voudra qu'elle}\unsure{,}} \inserted{r\unsure{é|e}sultera des axiomes} qu'elle ne dépend que de $\ll$. En tous cas, il est clair que la multiombre ne dépend que de $\underset{M}{\ll}$ (au lieu de $\underset{\mathcal{F}}{\ll}$).

\margin{\uline{NB} Je \unsure{vais|vois} changer plus bas le point de vue \unsure{pris|pris pour|pour} $\ll$ et les axiomes C4, C5.}

Je veux que cette multiombre soit une \uline{figure ensembliste dans} $M$, \uline{ayant pour support l'ombre}, \inserted{et indexée fid\unsure{è|e}lement par $\widetilde F$\,$\unsure{\phi}$} Cela nous amène à
\margin{$\phi$ avec la relation d'ordre $\leq$ dans $\widetilde F$ correspondant à l'inclusion de $\mathrm{Multomb}(F)$, i.e. $\downarrow$}

\inserted{\uline{C4} Si $X,Y\in M$ tels que $X\between Y$ (i.e. $X,Y$ font partie d'une = figure $F$), alors $X\ll Y\Rightarrow X\leq Y$}

\uline{C\deleted{\unsure{4}}5}: Si $X\in M$, $F\in\mathcal{F}$, \inserted{tels que $X\ll F$, et} si $F^X$ est l'ens
\[
 \text{\deleted{\unsure{des|plus}} des }Y\in M\quad\text{tels que }\deleted{\unsure{X\ll Y|X\leq Y}\,\unsure{:}}\,X\ll Y\leq F
 \tag*{(1.37)}
\]
(l'ens des strates de $F$ \unsure{raffinées|raffinée} par $X$), \deleted{\unsure{on}} cet ens $F^X$ . admet un plus petit élément \inserted{pour $\leq$} (i.e. il existe une strate de $F$ raffinée par $X$, et parmi \deleted{\unsure{toutes}} les strates de $F$ raffinées par $X$, il en est une plus petite)

\margin{\uline{Cor} \inserted{Si $X\in M$, $F\in\mathcal{F}$} Pour que $X\ll F$, il f et S qu'il existe $Y\deleted{\in\unsure{\widetilde F|\mathcal{F}}}\in M$ avec $\unsure{X|x}\ll\underbrace{Y\leq F}_{Y\in\widetilde F}$. Ainsi $X\ll F$ \unsure{dans|de} $M\times\mathcal{F}$ connu via $X\ll Y$ \unsure{dans|de} $M\times M$}

\deleted{\uline{C6} Soit $\unsure{F\in\mathcal{F},\,X\in\widetilde F}$ \unsure{puis|soit} $X\in M$, $\unsure{X\ll F|X\in\mathcal{F}}$.}

\deleted{Alors $\exists Y\in\unsure{M|\mathcal{F}M}$ avec \unsure{*|t.q.} a) $Y\ll X$ i.e. $\unsure{X|Y}\in F^Y$ \quad b) $Y\ll X'\leq X$ et $X'\in M\Longrightarrow X'\unsure{=|\simeq}X$}

\deleted{b) \unsure{que|pour} $X$ est \unsure{la|le} plus \uline{petit} \uline{\unsure{élément}} de \unsure{$F$} (cf C4)}

\margin{\uline{NB} Le fait que dans $\mathrm{Multom}(X)$, les $\mathrm{Omb}(X)^{\unsure{\circ|o}}\ne\varnothing$, est trivial\,; car $X\in\mathrm{Omb}(X)^\circ$}
